% Kapitel 9, satt ur lectures/L09/README.md.

\chapter{Praktisk tentamen 1: sekvensnät}
\label{c9:ch}

{\large\itshape Tre timmar, individuellt, vid datorn. Omfattar del~\ref{part:vhdl} i sin helhet, med
tyngdpunkt på sekvensnät.\par}

\begin{chapterbox}{I det här kapitlet}
\begin{itemize}
  \item Tentamens form: fristående uppgifter, var och en med en utdelad självkontrollerande
    testbänk.
  \item Vad som är tillåtet att använda, och vad som inte är det.
  \item Hur man förbereder sig, och vad som är den bästa förberedelsen som finns.
  \item Vad tentamen faktiskt mäter.
\end{itemize}
\end{chapterbox}

Det här är kursens första examinationstillfälle. Det innehåller ingen ny teori; kapitlet säger vad
tentamen mäter, vilken form uppgifterna har och hur man förbereder sig, och det är värt att läsa i
förväg snarare än dagen innan. Poäng och betygsgränser står i bilaga~\ref{app:project}.

Tentamen omfattar \chapref{c1:ch} till \ref{c8:ch} i sin helhet, med tyngdpunkt på sekvensnät:
register, räknare, skiftregister, timers, synkroniserare och enkla Mooremaskiner.

% ------------------------------------------------------------------------------------------
\section{Upplägg}
\label{c9:sec:format}

Tentamen består av ett antal fristående uppgifter. Varje uppgift ger dig en \textbf{utdelad,
självkontrollerande testbänk} och en beskrivning av den modul den driver: portordning, typer och
beteende. Din uppgift är att skriva modulen så att testbänken passerar.

\begin{itemize}
  \item Uppgifterna är oberoende av varandra. En modul du inte får att fungera stoppar ingen annan.
  \item Uppgifterna är ordnade ungefär efter stigande svårighet, men läs igenom alla innan du
    börjar.
  \item En modul som inte är helt klar rättas på det som finns. Delvis korrekt logik ger delpoäng,
    så lämna aldrig in en tom fil för att den inte blev färdig.
  \item Testbänken är facit under skrivningen, inte efter den. \code{all checks passed!} betyder att
    den uppgiften är klar; rättningen läser dessutom koden.
\end{itemize}

Ett övningsprov med uppgifter av samma slag och samma omfattning delas ut i förväg.

% ------------------------------------------------------------------------------------------
\section{Hjälpmedel}
\label{c9:sec:aids}

\textbf{Tillåtet:}
\begin{itemize}
  \item Allt kursmaterial: den här boken, föreläsningarnas README, samtliga appendix, och all VHDL
    du själv skrivit under del~\ref{part:vhdl}, inklusive dina lösningar på övningarna.
  \item Din egen dator, med GHDL. Se \secref{c2:sec:testbenches} om du behöver påminna dig om hur en
    testbänk körs.
  \item CircuitVerse, om du vill rita en krets innan du skriver den. Det är ofta snabbare, inte
    långsammare.
\end{itemize}

\textbf{Inte tillåtet:} kommunikation med andra i någon form; språkmodeller och andra verktyg som
genererar eller förklarar VHDL åt dig; och sökning på nätet efter färdiga lösningar. Gränsen går
inte vid vem som trycker på tangenterna utan vid vems förståelse som prövas. Att slå upp GHDL:s
eller VHDL:s dokumentation är tillåtet; att söka efter en färdig modul är det inte.

Den fullständiga och auktoritativa hjälpmedelslistan står i bilaga~\ref{app:project}; står något
annat här gäller den.

% ------------------------------------------------------------------------------------------
\section{Före tentamen}
\label{c9:sec:prep}

\begin{itemize}
  \item Gå igenom övningsprovet, under tidspress och utan att titta i lösningarna först.
  \item Övning~\ref{c8:ex:rx} och \ref{c8:ex:tx}, bokens två capstones, är den bästa förberedelsen
    som finns: de tvingar dig att konstruera ett tillståndsdiagram själv och komponera tidigare
    moduler runt det, vilket är precis vad tentamen mäter.
  \item Se till att GHDL fungerar på din maskin \emph{innan} du kommer. En trasig verktygskedja är
    inget skäl till förlängd skrivtid.
\end{itemize}

% ------------------------------------------------------------------------------------------
\section{Det här mäter tentamen}
\label{c9:sec:measures}

\begin{itemize}
  \item Att du kan skriva den synkrona processmallen korrekt, utan att slå upp den.
  \item Att du kan bygga en räknare, en timer och ett skiftregister och veta vad som skiljer dem.
  \item Att du synkroniserar varje asynkron ingång, utan att bli påmind.
  \item Att du kan konstruera en liten Mooremaskin ur en specifikation i ord, och skriva den som en
    uppräknad typ och ett \code{case}.
  \item Att du kan läsa en testbänks felutskrift och hitta tillbaka till raden som orsakade den.
\end{itemize}

% ------------------------------------------------------------------------------------------
\section{Efter tentamen}
\label{c9:sec:after}

Nästa kapitel startar \textbf{grupprojektet}, som löper över \chapref{c10:ch} till \ref{c19:ch} och
står för halva kursens poäng. Läs projektspecifikationen i bilaga~\ref{app:project} före
\chapref{c10:ch}; grupperna sätts samman där.

\Chapref{c10:ch} tar upp projektstarten: CAN-bussen, ramen och arbitreringen, och projektets första
VHDL, det delade paketet \code{can_def}.
